%============================================================
% MTH 112Z Project - Supplement Graphing Sinusoidal Functions
% Updated 202302
%============================================================

\section{Graphing Sinusoidal Functions} \label{supplement-graphing-sinusoidal-functions}
In this Section, we will learn some things.

  \section*{Graphing Sinusoidal Functions: Phase Shift vs. Horizontal Shift}
    
      Let's consider the function $g(x)=\sin \left( 2x-\frac{2\pi}{3} \right)$.  Using what we study in MTH 111 about graph transformations, it should be apparent that the graph of $g(x)=\sin \left( 2x-\frac{2\pi}{3} \right)$ can be obtained by transforming the graph of $g(x)=\sin(x)$. (To confirm this, notice that $g(x)$ can be expressed in terms of $f(x)=\sin(x),$ as $g(x)=f \left( 2x-\frac{2\pi}{3} \right)$.)  Since the constants $2$ and $\frac{2\pi}{3}$ are multiplied by and subtracted from the input variable, $x$, what we study in MTH 111 tells us that these constants represent a horizontal stretch/compression and a horizontal shift, respectively.
    
      It is often recommended in MTH 111 that we factor-out the horizontal stretching/compressing factor before transforming the graph, i.e., it's often recommended that we first re-write $g(x)=\sin \left( 2x-\frac{2\pi}{3} \right)$  as $g(x)=\sin \left( 2 \left( x-\frac{\pi}{3} \right)\right)$.
    
      After writing $g$ in this format, we can draw its graph by performing the following sequence of transformations of the base function $f(x)=\sin(x)$:
      
      \begin{enumerate}
        \item Compress horizontally by a factor of $\frac{1}{2}$.
        \item Shift $\frac{\pi}{3}$ units to the right.
      \end{enumerate}
    
      The advantage of this method is that the $y$-intercept of $f(x)=\sin(x)$, $ \left( 0,0 \right)$, ends-up exactly where the horizontal shift suggests: when we compress the graph by a factor of $\frac{1}{2}$, the the $y$-intercept of the graph doesn't move since $\frac{1}{2} \cdot 0 = 0$; then, when we shift the graph $\frac{\pi}{3}$ units to the right, the point  $\left( 0,0\right)$  ends up at $\left( \frac{\pi}{3},0\right)$; so the $y$-intercept ends up moving $\frac{\pi}{3}$ units to the right, exactly how far we shifted.
    
      Compare this with the alternative method: we can leave $g(x)=\sin  \left( 2x-\frac{2\pi}{3} \right)$ as-is and skip factoring-out the horizontal stretching/compressing factor, but then we need the following sequence to transform $f(x)=\sin(x)$ into the graph of $g$:
      
      \begin{enumerate}
        \item Shift $\frac{2\pi}{3}$ units to the right.
        \item Compress horizontally by a factor of $\frac{1}{2}$.
      \end{enumerate}
    
      The disadvantage of this method is that the $y$-intercept of $f(x)=\sin(x)$ \textbf{doesn't} end-up where the horizontal shift suggests: when we shift the graph of $f(x)=\sin(x)$ to the right by $\frac{2\pi}{3}$ units, the $y$-intercept moves from $ \left( 0,0 \right)$ to $\left( \frac{2\pi}{3},0\right)$; then, when we compress the graph by a factor of $\frac{1}{2}$, it moves to $\left( \frac{\pi}{3},0\right)$, so the $y$-intercept \textbf{doesn't} end up shifted $\frac{2\pi}{3}$ units to the right.
    
      Figure \ref{fig:shifts} shows the graphs of $y=f(x)$ and $y=g(x)$. Notice that the behavior of $y=g(x)$ at $x=\frac{\pi}{3}$ is like the behavior of $y=f(x)$ at $x=0$, i.e., $y=g(x)$ appears to have been shifted $\frac{\pi}{3}$ units to the right. For this reason, $\frac{\pi}{3}$ is called the horizontal shift of $g(x)=\sin  \left( 2x-\frac{2\pi}{3}\right)=\sin \left( 2\left( x-\frac{\pi}{3}\right)\right)$.
      
      \begin{center}
        \captionsetup{type=figure} \captionof{figure}{$y=g(x)$ with $f(x)=\sin(x)$}
          %<description>This is a graph showing both f(x) and g(x) where the function g(x) is a horizontal shift of f(x).</description>
            \begin{tikzpicture}
              \begin{axis}[
                  width=9cm,
                  height=6cm,
                  xmin=,xmax=,
                  ymin=-1.6, ymax=2.2,
                  xtick={-2*pi/3,-pi/3,...,8*pi/3},
                  xticklabels={$-\frac{2\pi}{3}$, $-\frac{\pi}{3}$, $0$, $\frac{\pi}{3}$, $\frac{2\pi}{3}$, $\pi$, $\frac{4\pi}{3}$, $\frac{5\pi}{3}$, $2\pi$, $\frac{7\pi}{3}$},
                  minor xtick = {-2*pi/3,-pi/2,...,3*pi},
                  ytick={-1,1},
                  minor ytick={-1.5,-1,...,2},
                  clip=false,
                ]
                \addplot[trigfunction, black, domain=-2.5:8.5, samples=100] {sin(deg(x))};
                \addplot[trigfunction, domain=-2.5:8.5, samples=100] {sin(deg(2*x)-120))};
                \legend{$y=f(x)$,$y=g(x)$}
              \end{axis}
            \end{tikzpicture}\label{fig:shifts}
      \end{center}
    
      The constant $\frac{2\pi}{3}$ is given a different name, phase shift, since it can be used to determine how far out-of-phase a sinusoidal function is in comparison with $y=\sin(x)$ or $y=\cos(x)$. To determine how far out-of-phase a sinusoidal function is, we can determine the ratio of the phase shift and $2\pi$. (We use $2\pi$ because it's the period of $y=\sin(x)$ and $y=\cos(x)$.)  Since $\frac{2\pi}{3}$ is the phase shift for $g(x)=\sin \left( 2x-\frac{2\pi}{3} \right)$, the graph of $y=g(x)$ is out-of-phase $\frac{2\pi/3}{2\pi}=\frac{1}{3}$ of a period. (Since this number is positive, it represents a horizontal shift to the right $\frac{1}{3}$ of a period.)
    

    \begin{myDefinition}
          Given a sinusoidal function of the form $y=A\sin(wx-C)+k$ or $y=A\cos(wx-C)+k$, the \textbf{phase shift} is $C$ and $\frac{|C|}{2\pi}$ represents the fraction of a period that the graph has been shifted (shift to the right if $C$ is positive or to the left if $C$ is negative).
    \end{myDefinition}
    \begin{myDefinition}
          If we re-write the function as $y=A\sin \left( w \left( x-\frac{C}{w}\right)\right)+k$ or $y=A\cos  \left( w \left( x-\frac{C}{w}\right)\right)+k$, we can see that the \textbf{horizontal shift} is $\frac{C}{w}$ units (shift to the right if $\frac{C}{w}$ is positive or to the left if $\frac{C}{w}$ is negative).
    \end{myDefinition}
    <example>
        
          Identify the phase shift and horizontal shift of $g(x)=\cos \left( 3x-\frac{\pi}{4} \right)$.
        
      <solution>
        
          The phase shift of $g(x)=\cos \left( 3x-\frac{\pi}{4} \right)$ is $\frac{\pi}{4}$. This tells us that the graph of $y=g(x)$ is out of phase $\frac{|\pi/4|}{2\pi}=\frac{1}{8}$ of a period, i.e., compared with $y=\cos(x)$, the graph of $g(x)=\cos(3x-\frac{\pi}{4})$ has been shifted one-eighth of a period to the right.
        
        
          To find the horizontal shift, we need to factor-out $3$ from $3x-\frac{\pi}{4}$.
            \begin{align*}
              g(x) &=\cos\left( 3x-\frac{\pi}{4}\right)\\
               &=\cos\left(3\left(x-\frac{\pi}{3 \cdot 4}\right)\right)\\
               &=\cos\left(3\left(x-\frac{\pi}{12}\right)\right)\\
            \end{align*}
        
        
          So the horizontal shift is $\frac{\pi}{12}$. This tells us, that compared with $y=\cos(x)$, the graph of $g(x)=\cos(3x-\frac{\pi}{4})$ has been shifted $\frac{\pi}{12}$ to the right.
        
        
          Notice that the period of $g(x)=\cos(3x-\frac{\pi}{4})$ is $2\pi \cdot \frac{1}{3}=\frac{2\pi}{3}$, and one-eighth of $\frac{2\pi}{3}$ is $\frac{2\pi}{3} \cdot \frac{1}{8}=\frac{\pi}{12}$, so a shift of one-eighth of a period is the same as a shift of $\frac{\pi}{12}$ units!
        
      </solution>
    </example>
    <example>
      <statement>
        
          Draw a graph $q(t)=2\sin(4t+\pi)+1$. First, find it's amplitude, period, midline, phase shift, and horizontal shift.
        
      </statement>
      <solution>
        
          Amplitude: $|A|=|2|=2$
          Period: $P=2\pi \cdot \frac{1}{|w|}=\frac{2\pi}{4}=\frac{\pi}{2}$
          Midline: $y=1$
          Phase shift: $-\pi$ (this tells is that the graph is out-of-phase $\frac{|-\pi|}{2\pi}=\frac{1}{2}$ of a period)
        
        
          Horizontal shift: $\frac{\pi}{4}$ units to the left since:
            \begin{align*}
              q(t) &=2\sin\left( 4t+\pi \right)+1\\
               &=2\sin \left(4 \left(t+\frac{\pi}{4} \right)\right)+1\\
               &=2\sin\left( 4 \left(t- \left( -\frac{\pi}{4}\right)\right)\right)+1\\
            \end{align*}
        
          Now we can draw a graph of $q(t)=2\sin(4t+\pi)+1$ by drawing a sinusoidal function with the necessary features; seeFigure \ref{fig:sinusoidal}.
        
        \begin{center}
          \captionsetup{type=figure} \captionof{figure}{$y=q(t)$}
            %<description>Graph of q(t) = 2sin(4t+pi)+1.</description>
              \begin{tikzpicture}
                \begin{axis}[
                    width=9cm,
                    height=6cm,
                    xmin=,xmax=,
                    ymin=-2.5, ymax=4.5,
                    xtick={-pi/4,0,...,pi},
                    xticklabels={$-\frac{\pi}{4}$, $0$, $\frac{\pi}{4}$, $\frac{\pi}{2}$, $\frac{3\pi}{4}$, $\pi$},
                    minor xtick = {-pi/4,-pi/8,...,pi},
                    ytick={-2,-1,...,4},
                    xlabel={$t$},
                    clip=false,
                 ]
                  \addplot[trigfunction, domain=-1:3.5, samples=100] {2*sin(deg(4*x)+180)+1};
                \end{axis}
              \end{tikzpicture}\label{fig:sinusoidal}
        \end{center}

      </solution>
    </example>

  <exercises xml:id="exercises-sinusoidal-functions">
    <exercisegroup cols="2" xml:id="exercisegroup-sketch-waves">

      \subsection*{Sketch Sinusoidal Graphs}
          In <xref ref="exercisegroup-sketch-waves" text="local">Exercises</xref>, draw a graph of each of the following functions. List the amplitude, midline, period, phase shift, and horizontal shift.
      <exercise>
        <statement>
            $f(x)=3\sin \left( 3x-\frac{\pi}{2} \right)$
        </statement>
        <answer>
            \begin{enumerate}
              \item Amplitude $3$ units
              \item Period $\frac{2\pi}{3}$ units
              \item Midline $y=0$
              \item Phase shift $\frac{\pi}{2}$
              \item Horizontal Shift $\frac{\pi}{6}$ units to the right
            \end{enumerate}
          
            %<description>The graph of f(x)=3sin(3x-pi/2).</description>
            
              \begin{tikzpicture}
                \begin{axis}[
                    width=6cm,
                    xmin=,xmax=,
                    ymin=-4, ymax=4,
                    xtick={-pi/3,-pi/6,...,pi},
                    xticklabels={$-\frac{\pi}{3}$, $-\frac{\pi}{6}$, $0$, $\frac{\pi}{6}$, $\frac{\pi}{3}$, $\frac{\pi}{2}$, $\frac{2\pi}{3}$, $\frac{5\pi}{6}$, $\pi$},
                    ytick={-3,3},
                    minor ytick={-3,-2,...,3},
                    clip=false,
                  ]
                  \addplot[trigfunction, domain=-1.5:3.5, samples=150] {3*sin(deg(3*x)-90)};
                  \addplot[soliddot] coordinates {(pi/6,0) (pi/3,3) (pi/2,0) (2*pi/3,-3) (5*pi/6,0)};
                \end{axis}
              \end{tikzpicture}
            
        </answer>
      </exercise>
      <exercise>
        <statement>
          
            $g(t)=\cos(4t+\pi)+3$
          
        </statement>
        <answer>
            \begin{enumerate}
              \item Amplitude $1$ unit 
              \item Period $\frac{\pi}{2} $ units 
              \item Midline $y=3$ 
              \item Phase shift $-\pi$ 
              \item Horizontal Shift $\frac{\pi}{4}$ units to the left 
            \end{enumerate}
          
            %<description>The graph of g(t)=cos(4t+pi)+3.</description>
            
              \begin{tikzpicture}
                \begin{axis}[
                    width=6cm,
                    xmin=,xmax=,
                    xlabel={$t$},
                    ymin=-1.2, ymax=4.5,
                    xtick={-pi/2,-pi/4,...,5*pi/4},
                    xticklabels={$-\frac{\pi}{2}$, $-\frac{\pi}{4}$, $0$, $\frac{\pi}{4}$, $\frac{\pi}{2}$, $\frac{3\pi}{4}$, $\pi$, $\frac{5\pi}{4}$},
                    ytick={1,2,3,4},
                    clip=false,
                  ]
                  \addplot[trigfunction, domain=-1.8:4, samples=150] {cos(deg(4*x)+180)+3};
                  \addplot[guideline] coordinates {(-1.7,3) (3.9,3)};
                  \addplot[soliddot] coordinates {(-pi/4,4) (-pi/8,3) (0,2) (pi/8,3) (pi/4,4)};
                \end{axis}
              \end{tikzpicture}
            
          
        </answer>
      </exercise>
      <exercise>
        <statement>
          
            $m(\theta)=2 \cos \left( 2\pi \theta - \pi \right)+4$
          
        </statement>
        <answer> 
              \item Amplitude $2$ units.
              \item Period $1$ unit.
              \item Midline $y=4$.
              \item Phase shift $\pi$.
              \item Horizontal Shift $\frac{1}{2}$ of a unit to the right.
            %<description>The graph of m(theta)=2cos(2 pi theta - pi) + 4.</description>
            
              \begin{tikzpicture}
                \begin{axis}[
                    width=6cm,
                    xmin=,xmax=,
                    ymin=-3, ymax=7,
                    xlabel={$\theta$},
                    xtick={-1,1,2},
                    minor xtick = {-1,-0.75,...,1.75},
                    ytick={-2,2,4,6},
                    minor ytick = {-2,-1,...,6},
                    clip=false,
                  ]
                  \addplot[trigfunction, domain=-1.2:2.2, samples=100] {2*cos(deg(2*pi*x)-180)+4};
                  \addplot[guideline] coordinates {(-1.2,4) (2.2,4)};
                  \addplot[soliddot] coordinates {(1/2,6) (3/4,4) (1,2) (5/4,4) (3/2,6)};
                \end{axis}
              \end{tikzpicture}
            
        </answer>
      </exercise>
      <exercise>
        <statement>
          
            $n(x)=-4\sin \left( \pi x+\frac{\pi}{4} \right)-2$
          
        </statement>
        <answer>
          
            <dl width="narrow">
              \item
                \subsection*{Amplitude}
                
                  $4$ units
                
              
              \item
                \subsection*{Period}
                
                  $2$ units
                
              
              \item
                \subsection*{Midline}
                
                  $y=-2$
                
              
              \item
                \subsection*{Phase shift}
                
                  $-\frac{\pi}{4}$
                
              
              \item
                \subsection*{Horizontal Shift}
                
                  $\frac{1}{4}$ of a unit to the left
                
              
            </dl>
          
          
            %<description>The graph of n(x)=-4sin(pi x + pi/4)-2.</description>
            
              \begin{tikzpicture}
                \begin{axis}[
                    width=6cm,
                    xmin=,xmax=,
                    ymin=-8, ymax=4,
                    xtick={-1,1,2,3},
                    minor xtick = {-1,-0.75,...,3.5},
                    ytick={-6,-4,-2,2},
                    minor ytick = {-7,-6,...,3},
                    clip=false,
                  ]
                  \addplot[trigfunction, domain=-1.25:3.75, samples=100] {-4*sin(deg(pi*x)+45)-2};
                  \addplot[guideline] coordinates {(-1.25,-2) (3.75,-2)};
                  \addplot[soliddot] coordinates {(-1/4,-2) (1/4,-6) (3/4,-2) (5/4,2) (7/4,-2)};
                \end{axis}
              \end{tikzpicture}
            
          
        </answer>
      </exercise>
    </exercisegroup>
    <exercisegroup cols="2" xml:id="exercisegroup-find-sine-formulas">
      \subsection*{Find the Formula}
      
        
          In <xref ref="exercisegroup-find-sine-formulas" text="local">Exercises</xref>, find two algebraic rules (one involving sine and one involving cosine) for each of the functions graphed below.
        
      
      <exercise>
        <statement>
          
            $y=p(t)$
          
            %<description>Graph of y=p(t). The highest points are at 2 and the lowest are at -6. On the y-axis the intercept is (0, -6). The period is pi.</description>
            
              \begin{tikzpicture}
                \begin{axis}[
                    width=6cm,
                    xmin=,xmax=,
                    xlabel={$t$},
                    ymin=-6.75, ymax=2.75,
                    xtick={-pi/2,-pi/4,...,5*pi/4},
                    xticklabels={$-\frac{\pi}{2}$, $-\frac{\pi}{4}$, $0$, $\frac{\pi}{4}$, $\frac{\pi}{2}$, $\frac{3\pi}{4}$, $\pi$, $\frac{5\pi}{4}$},
                    ytick={-6,-4,-2,2},
                    clip=false,
                  ]
                  \addplot[trigfunction, domain=-2.3:4.6, samples=100] {4*sin(deg(2*x)-90)-2};
                \end{axis}
              \end{tikzpicture}
            
          
        </statement>
        <answer>
          
            $ \begin{aligned} p(x)=4 \sin \left( 2 \left( x-\frac{\pi}{4} \right)\right)-2 \\ p(x)=4 \cos \left( 2 \left( x-\frac{\pi}{2} \right)\right)-2 \end{aligned} $
          
        </answer>
      </exercise>
      <exercise>
        <statement>
          
            $y=q(x)$
          
          
            %<description>Graph of y=q(t). The highest points are at 2 and the lowest are at -4. It has points at (0.25, 2), (0.75, -1), and (1.25, -4). The period is 2.</description>
            
              \begin{tikzpicture}
                \begin{axis}[
                    width=6cm,
                    xmin=,xmax=,
                    ymin=-5, ymax=3,
                    xtick={-2,-1,...,3},
                    minor xtick = {-2.25,-2,...,3.25},
                    ytick={-4,-2,...,2},
                    minor ytick = {-3,-1,1},
                    clip=false,
                    major grid style={thick},
                  ]
                  \addplot[trigfunction, domain=-2.5:3.5, samples=100] {3*sin(deg(pi*x)+45)-1};
                \end{axis}
              \end{tikzpicture}
            
          
        </statement>
        <answer>
          
              \begin{aligned}
                q(x)&=3 \sin \left( \pi \left( x+\frac{1}{4} \right)\right)-1\\
                q(x)&=3 \cos \left( \pi \left( x-\frac{1}{4} \right)\right)-1\\
              \end{aligned}
          
